% https://ctan.org/tex-archive/macros/latex/contrib/fp?lang=en
% https://mirrors.ibiblio.org/CTAN/macros/generic/xint/xint.pdf

\newcommand{\trans}[1]{\ensuremath{{#1}\mkern-1mu^{\mathsf{T}}}}
\newcommand{\norm}[1]{\ensuremath{\left\Vert #1 \right\Vert}}

\usepackage{xinttools} % for the \xintFor
\usepackage{xintfrac}
\usepackage{fp}
\usepackage{xstring}

%%%%%%%%%%%%%%%%%%%%%%%%%%%%%%%%%%%%%%%%%%%%%%%%%%%%%%%%%%%
% Print Utility
%%%%%%%%%%%%%%%%%%%%%%%%%%%%%%%%%%%%%%%%%%%%%%%%%%%%%%%%%%%
%%%%%%%%%%%
% Cleanup %
%%%%%%%%%%%
\newcommand\roundNumber[2][4]{%
    \FPset\printNum{#2}%
    \FPround\printNum\printNum{#1}
    \FPclip\printNum\printNum
    \printNum
}

%%%%%%%%%%%%%%%
% Conversions %
%%%%%%%%%%%%%%%
% Convert a fraction with root to a decimal number
% TODO make #3 optional
\newcommand\fractionToNum[3]{
    \FProot\tmpA{#3}{2}
    \FPmul\tmpA\tmpA{#2}
    \FPdiv\tmpB{#1}\tmpA
    \FPset\fractionNumRes{\tmpB}
    \ensuremath{
        \tmpB
    }
}

% Convert a decimal number to a fraction with root
\newcommand\solutionC{%
    \FPset\finalmul{1}
    \FPset\tmptmptmptmp\tmptmptmp
    %
    \edef\innersequence{\xintSeq[+1]{+1}{+10}}%
    \xintFor* ##1 in \innersequence \do {%
        \FPtrunc\floored\tmptmptmptmp{0}
        \FPsub\remainder\tmptmptmptmp\floored
        \FPifgt{\remainder}{0.00001}
            \FPdiv\mulmul{1}{\remainder}
            \FPmul\finalmul\finalmul\mulmul
        \else
            \xintBreakFor
        \fi
        \FPset\tmptmptmptmp\mulmul
    }
    \FPround\finalmul\finalmul{8}
    \FPclip\finalmul\finalmul
    \FPset\finalmultmp\finalmul
    % The result will be a multiplicator that is able to round
    % the number to a whole number, but the result might not be a
    % nice number to root e.g. 8/sqrt(6), 0.09375 will result in 32
    % 0.09375 * 32 = 3, but sqrt(32) = 5,65.. but given this we know
    % that the searched for number must be a multiple of 32, in this case 64
    % Note: rounding to 6 & 36 does not work since 0.09375 * 36 = 3.375
    \xintFor* ##1 in \innersequence \do {%
        \FProot\finalmulsqrt\finalmultmp{2}
        \FPround\finalmulsqrt\finalmulsqrt{4}
        \FPclip\finalmulsqrt\finalmulsqrt
        \FPifint{\finalmulsqrt}
            \xintBreakFor
        \else
            \FPadd\finalmultmp\finalmultmp\finalmul
        \fi
    }
    \FPset\finalmul\finalmultmp
    \FProot\finalmulsqrt\finalmul{2}
    \FPround\finalmulsqrt\finalmulsqrt{4}
    \FPclip\finalmulsqrt\finalmulsqrt
    \FPmul\divisor\divisor\finalmul
    \FPmul\dividend\dividend\finalmulsqrt
}

\newcommand\numToFraction[1]{%
    \FPset\dividend{1}
    \FPset\divisor{#1}
    \ensuremath{
        % reverse fraction
        \FPdiv\divisor{1}{#1}
        % Reverse root
        % Use FPmul instead of FPpow to allow for neg values
        \FPmul\divisor\divisor\divisor
        \FPset\tmptmptmp\divisor
        \FPround\divisor\divisor{10}
        \FPclip\divisor\divisor
        \FPifint{\divisor}
        \else
            \solutionC
        \fi
        %
        \FPround\divisor\divisor{4}
        \FPclip\divisor\divisor
        \FPround\dividend\dividend{4}
        \FPclip\dividend\dividend
        %
        \FPset\outputVar{#1}
        \FPset\outputVar{\frac{\dividend}{\sqrt{\divisor}}}
        \FPset\inputVar{#1}
        \FPmul\inputVar\inputVar\inputVar
        \FPround\inputVar\inputVar{8}
        \FPclip\inputVar\inputVar
        \FPifint{\inputVar}
            \FPset\outputVar{\sqrt{\inputVar}}
        \else
        \fi
        \FPset\inputVarA{#1}
        \FPround\inputVarA\inputVarA{10}
        \FPclip\inputVarA\inputVarA
        \FPifint{\inputVarA}
            \FPset\outputVar{\inputVarA}
        \else
        \fi
        \outputVar
    }
}

\newcommand\numToFractionA[1]{%
    \FPset\inputVarA{#1}
    \FPround\inputVarA\inputVarA{10}
    \FPclip\inputVarA\inputVarA
    \FPifint{\inputVarA}
        \FPset\outputVar{\inputVarA}
        \outputVar
    \else
        \numToFraction{#1}
    \fi
}

%%%%%%%%%%%%%%%%%%%%%%%%%%%%%%%%%%%%%%%%%%%%%%%%%%%%%%%%%%%
% Basic number calculations
%%%%%%%%%%%%%%%%%%%%%%%%%%%%%%%%%%%%%%%%%%%%%%%%%%%%%%%%%%%
% Divide skalar by skalar
\newcommand\divNumber[3][4] {%
    \FPset\divNumRes{0}
    \FPset\dividend{#2}
    \FPdiv\divNumRes\dividend{#3}
    % \FPround\divNumRes\divNumRes{#1}
    % \FPclip\divNumRes\divNumRes
}

%%%%%%%%%%%%%%%%%%%%%%%%%%%%%%%%%%%%%%%%%%%%%%%%%%%%%%%%%%%
% Vector Create & Print
%%%%%%%%%%%%%%%%%%%%%%%%%%%%%%%%%%%%%%%%%%%%%%%%%%%%%%%%%%%
% Create a vector
\newcommand\newVec[2]{% {Basename}{input}
    \expandafter\edef\csname #1\endcsname{\xintCSVtoList{#2}}%
}

% Append to a vector
\newcommand\appendToVec[3]{% {Basename}{input}
    \expandafter\edef\csname #1\endcsname{#2 \xintCSVtoList{#3}}%
}

\newcommand\printVecT[2][0]{%
    \ensuremath{\trans{\left(
        \xintFor* ##1 in #2 \do {%
            \FPifeq{#1}{0}
                \roundNumber{##1}
            \else
                \numToFractionA{##1}
            \fi
            \xintifForLast{}{%
                ,
            }
        }
    \right)}}
}

\newcommand\printVec[2][0]{%
    \ensuremath{%
        \begin{pmatrix}
            \xintFor* ##1 in #2 \do {%
                \FPifeq{#1}{0}
                    \roundNumber{##1}
                \else
                    \numToFractionA{##1}
                \fi
                \xintifForLast{}{%
                    \\
                }
            }
        \end{pmatrix}
    }
}



%%%%%%%%%%%%%%%%%%%%%%%%%%%%%%%%%%%%%%%%%%%%%%%%%%%%%%%%%%%
% Vector Math
%%%%%%%%%%%%%%%%%%%%%%%%%%%%%%%%%%%%%%%%%%%%%%%%%%%%%%%%%%%

\newcommand\printSubVec[2] {%
    \subVec{#1}{#2}
    \printVec{\subVecResult}
}

% Subract list of vectors from first vector in list
\newcommand\subVec[2][subVecResult] {%
    \def\veclist{\xintCSVtoList{#2}}
    \edef\elemOne{\xintNthElt{1}{\veclist}}
    \edef\innersequence{\xintSeq[+1]{+1}{\expandafter\xintLength\expandafter{\elemOne}}}%
    \xintFor* ##1 in \innersequence \do {%
        \FPset\subCoord{0}%
        \xintifForFirst{%
            \xintFor* ##2 in \veclist \do {%
                \xintifForFirst{%
                    \FPset\subCoord{\xintNthElt{##1}{##2}}%
                }{%
                    \FPsub\subCoord\subCoord{\xintNthElt{##1}{##2}}%
                }
            }
            \newVec{#1}{\subCoord}
        }{%
            \xintFor* ##2 in \veclist \do {%
                \xintifForFirst{%
                    \FPset\subCoord{\xintNthElt{##1}{##2}}%
                }{%
                    \FPsub\subCoord\subCoord{\xintNthElt{##1}{##2}}%
                }
            }
            \appendToVec{#1}{\csname #1\endcsname}{\subCoord}
        }
    }
}

% Sum list of vectors
\newcommand\sumVec[2][sumVecResult] {%
    \def\veclist{\xintCSVtoList{#2}}
    \edef\elemOne{\xintNthElt{1}{\veclist}}
    \edef\innersequence{\xintSeq[+1]{+1}{\expandafter\xintLength\expandafter{\elemOne}}}%
    \xintFor* ##1 in \innersequence \do {%
        \FPset\sumCoord{0}%
        \xintifForFirst{%
            \xintFor* ##2 in \veclist \do {%
                \FPadd\sumCoord\sumCoord{\xintNthElt{##1}{##2}}%
            }
            \newVec{#1}{\sumCoord}
        }{%
            \xintFor* ##2 in \veclist \do {%
                \FPadd\sumCoord\sumCoord{\xintNthElt{##1}{##2}}%
            }
            \appendToVec{#1}{\csname #1\endcsname}{\sumCoord}
        }
    }
}

% Divide vector by skalar
\newcommand\divVec[3][divVecResult] {%
    \FPifeq{#3}{0}
        \PackageError{math package}{In div vector by scalar}{divisor is 0}
    \else
    \fi
    \xintFor* ##1 in #2 \do {%
        \FPdiv\divNumRes{##1}{#3}
        \xintifForFirst{%
            \newVec{#1}{\divNumRes}
        }{%
            \appendToVec{#1}{\csname #1\endcsname}{\divNumRes}
        }
    }
}

% Multiply vector by skalar
\newcommand\mulVec[3][mulVecResult] {%
    \xintFor* ##1 in #2 \do {%
        \FPmul\mulNumRes{##1}{#3}
        \xintifForFirst{%
            \newVec{#1}{\mulNumRes}
        }{%
            \appendToVec{#1}{\csname #1\endcsname}{\mulNumRes}
        }
    }
}

%%%%%%%%%%%%%%%%%%
% Norm of vector %
%%%%%%%%%%%%%%%%%%
\newcommand\printNorm[2][2] {%
    \FPset\tmp{0}
    \calcNorm{#2}
    \FPround\tmp\normResult{#1}
    \FPclip\tmp\tmp
    \tmp
}

\newcommand\printNormSqrt[1] {%
    \FPset\tmp{0}
    \calcNorm{#1}
    \numToFractionA{\normResult}
}

\newcommand\calcNorm[1] {%
    \FPset\normTotal{0}
    \FPset\iterVal{0}
    % Iterate each value, multiply it with itself and sum up
    \xintFor* ##1 in #1 \do {%
        % No idea why this needs to be abs,
        % but negative values are a problem here
        \FPabs\iterVal{##1}
        \FPpow\tmpNormB\iterVal{2}
        \FPadd\normTotal\normTotal\tmpNormB
    }
    % \FPround\normTotal\normTotal{4}
    % \FPclip\normTotal\normTotal
    % Evaluate root for numbered result
    \FProot\normRooted\normTotal{2}
    % \FPround\normRooted\normRooted{4}
    % \FPclip\normRooted\normRooted
    % Set output var
    \def\normResult{}
    \FPset\normResult{\normRooted}
}

%%%%%%%%%%%%%%%%%%%%%%%%%%%%%%
% Skalarproduct / Dotproduct %
%%%%%%%%%%%%%%%%%%%%%%%%%%%%%%
\newcommand\printSkalarp[3][2] {%
    \FPset\tmp{0}
    \calcSkalarp{#2}{#3}
    \FPround\tmp\scalarPResult{#1}
    \FPclip\tmp\tmp
    \tmp
}

\newcommand\calcSkalarp[2] {%
    \FPset\res{0}
    \FPset\tmp{0}
    \xintFor ##1 in {1,2,3} \do {%
        \FPset\valueA{\xintNthElt{##1}{#1}}
        \FPset\valueB{\xintNthElt{##1}{#2}}
        \FPmul\tmp\valueA\valueB
        \FPadd\res\res\tmp
    }
    \def\scalarPResult{}
    \FPset\scalarPResult{\res}
}

%%%%%%%%%%%%%%%%%%%%%%%%%%%%%%%%%%%%%%%%%%%%%%%%%%%%%%%%%%%
% Matrix Create & Print
%%%%%%%%%%%%%%%%%%%%%%%%%%%%%%%%%%%%%%%%%%%%%%%%%%%%%%%%%%%
